% Distilled blueprint for Evans, Chapter 10 (Hamilton--Jacobi Equations).
% Formal declarations, metadata, and citation identifiers are preserved; exposition and exercises are omitted.

\section{Introduction, Viscosity Solutions}
\label{sec:hj-introduction-viscosity-solutions}

Let $H:\R^n\times\R^n\to\R$ and $g:\R^n\to\R$ be continuous, and consider
\[
\begin{cases}
u_t+H(Du,x)=0 & \text{in }\R^n\times(0,\infty),\\
u=g & \text{on }\R^n\times\{t=0\}.
\end{cases}
\tag{1}
\]
The vanishing-viscosity approximation is
\[
\begin{cases}
u_t^\epsilon+H(Du^\epsilon,x)-\epsilon\Delta u^\epsilon=0
& \text{in }\R^n\times(0,\infty),\\
u^\epsilon=g & \text{on }\R^n\times\{t=0\},
\end{cases}
\tag{2}
\]
and suppose along a sequence $\epsilon_j\downarrow0$ that
\[
u^{\epsilon_j}\longrightarrow u
\quad\text{locally uniformly in }\R^n\times[0,\infty).
\tag{3}
\]

\subsection{Definitions}

\begin{remark}[Motivation for definition of viscosity solution]
\label{rem:hj-viscosity-solution-motivation}
\dcref{ch10:10.1-rem-1}
\group{Viscosity Solutions}
Fix $v\in C^\infty(\R^n\times(0,\infty))$ and suppose
\[
u-v\text{ has a strict local maximum at }(x_0,t_0).
\tag{5}
\]
Local uniform convergence (3) yields points $(x_{\epsilon_j},t_{\epsilon_j})\to(x_0,t_0)$ at which $u^{\epsilon_j}-v$ has a local maximum. At these points,
\[
Du^{\epsilon_j}=Dv,
\qquad
u_t^{\epsilon_j}=v_t,
\qquad
\Delta u^{\epsilon_j}\leq\Delta v.
\]
Equation (2) therefore gives
\[
v_t(x_{\epsilon_j},t_{\epsilon_j})
+H(Dv(x_{\epsilon_j},t_{\epsilon_j}),x_{\epsilon_j})
\leq\epsilon_j\Delta v(x_{\epsilon_j},t_{\epsilon_j}),
\tag{11}
\]
and passage to the limit gives
\[
v_t(x_0,t_0)+H(Dv(x_0,t_0),x_0)\leq0.
\tag{12}
\]
A quadratic perturbation of $v$ reduces a non-strict local maximum to (5). Reversing the argument at a local minimum gives
\[
v_t(x_0,t_0)+H(Dv(x_0,t_0),x_0)\geq0.
\tag{14}
\]
\end{remark}

\begin{definition}[Viscosity solution of the initial-value problem]
\label{def:viscosity-solution-cauchy-problem}
\uses{rem:hj-viscosity-solution-motivation}
\dcref{ch10:10.1-def-1}
\group{Viscosity Solutions}

A bounded, uniformly continuous function $u$ is called a \textit{viscosity solution} of the \textit{initial-value problem} (1) for the \textit{Hamilton-Jacobi equation} provided:

(i) $u = g$ on $\R^n \times \{t = 0\}$,

and

(ii) for each $v \in C^\infty(\R^n \times (0,\infty))$,

if $u - v$ has a local maximum at a point $(x_0, t_0) \in \R^n \times (0,\infty)$, then
\[
v_t(x_0, t_0) + H(Dv(x_0, t_0), x_0) \leq 0,
\tag{16}
\]
and

if $u - v$ has a local minimum at a point $(x_0, t_0) \in \R^n \times (0,\infty)$, then
\[
v_t(x_0, t_0) + H(Dv(x_0, t_0), x_0) \geq 0.
\tag{17}
\]
\end{definition}

\begin{remark}[The defining inequalities, not the vanishing-viscosity construction, are the definition]
\label{rem:viscosity-solution-definition-emphasis}
\uses{def:viscosity-solution-cauchy-problem}
\dcref{ch10:10.1-rem-2}
\group{Viscosity Solutions}

The intrinsic definition is the pair of test-function inequalities (16), (17). Vanishing viscosity is one construction that implies these inequalities, not part of the definition.
\end{remark}

\subsection{Consistency}

\begin{proposition}[Classical solutions are viscosity solutions]
\label{prop:classical-solutions-are-viscosity-solutions}
\uses{def:viscosity-solution-cauchy-problem}
\dcref{ch10:10.1-prop-1}
\group{Viscosity Solutions}

If $u \in C^1(\R^n \times [0,\infty))$ solves (1) and if $u$ is bounded and uniformly continuous, then $u$ is a viscosity solution. That is, any \textit{classical solution} of $u_t + H(Du,x) = 0$ \textit{is also a viscosity solution}.
\end{proposition}

\begin{proof}
\sketch At a local maximum or minimum of $u-v$, differentiability gives $Du=Dv$ and $u_t=v_t$. Substitution into the classical Hamilton--Jacobi equation gives equality, hence both viscosity inequalities.
\end{proof}

\begin{lemma}[Touching by a $C^1$ function]
\label{lem:touching-by-c1-function}
\dcref{ch10:10.1-lem-1}
\group{Viscosity Solutions}
Assume $u : \R^n \to \R$ is continuous and is also differentiable at some point $x_0$. Then there exists a function $v \in C^1(\R^n)$ such that
\[
u(x_0) = v(x_0)
\tag{18}
\]
and
\[
u - v \text{ has a strict local maximum at } x_0.
\tag{19}
\]
\end{lemma}

\begin{proof}
\sketch Subtract the affine tangent so that $x_0=0$ and $u(0)=Du(0)=0$. Write $u(x)=|x|\rho_1(x)$, dominate $|\rho_1|$ by a continuous nondecreasing radial modulus $\rho_2$, and set $v(x)=\int_{|x|}^{2|x|}\rho_2(r)\,dr+|x|^2$. Then $v\in C^1$, $u(0)=v(0)$, and $u(x)-v(x)\leq-|x|^2$ near zero.
\end{proof}

\begin{theorem}[Consistency of viscosity solutions]
\label{thm:consistency-viscosity-solutions}
\uses{def:viscosity-solution-cauchy-problem, lem:touching-by-c1-function}
\dcref{ch10:10.1-thm-1}
\group{Viscosity Solutions}

Let $u$ be a viscosity solution of (1), and suppose $u$ is differentiable at some point $(x_0, t_0) \in \R^n \times (0, \infty)$. Then
\[
u_t(x_0, t_0) + H(Du(x_0, t_0), x_0) = 0.
\]
\end{theorem}

\begin{proof}
\sketch The touching lemma in space-time supplies a $C^1$ function touching $u$ strictly from above at the differentiability point. Mollify it; nearby maxima converge to the contact point, so the viscosity subsolution inequality passes to the limit and gives the PDE inequality $\leq0$. Apply the same argument to $-u$ for the reverse inequality.
\end{proof}

\section{Uniqueness}
\label{sec:hj-uniqueness}

Fix $T>0$ and consider
\[
\begin{cases}
u_t+H(Du,x)=0 & \text{in }\R^n\times(0,T],\\
u=g & \text{on }\R^n\times\{t=0\}.
\end{cases}
\tag{1}
\]

\begin{definition}[Viscosity solution on a finite time interval]
\label{def:viscosity-solution-finite-horizon}
\uses{def:viscosity-solution-cauchy-problem}
\dcref{ch10:10.2-def-1}
\group{Viscosity Solutions}

A bounded, uniformly continuous function $u$ is a viscosity solution of (1) if $u=g$ at $t=0$ and, for each smooth test function $v$, a local maximum of $u-v$ in $\R^n\times(0,T)$ implies (16), while a local minimum implies (17).
\end{definition}

\begin{lemma}[Extrema at a terminal time]
\label{lem:extrema-at-terminal-time}
\uses{def:viscosity-solution-finite-horizon}
\dcref{ch10:10.2-lem-1}
\group{Viscosity Solutions}

Assume $u$ is a viscosity solution of (1) and $u - v$ has a local maximum (minimum) at a point $(x_0, t_0) \in \R^n \times (0, T]$. Then
\[
v_t(x_0, t_0) + H(Dv(x_0, t_0), x_0) \le 0 \quad (\ge 0).
\tag{2}
\]
Thus the viscosity inequalities extend to extrema with $t_0=T$.
\end{lemma}

\begin{proof}
\sketch For a strict maximum at $(x_0,T)$, replace $v$ by $v+\epsilon/(T-t)$. The perturbed difference attains an interior maximum converging to $(x_0,T)$; apply the viscosity inequality and let $\epsilon\downarrow0$. The minimum case uses the opposite barrier.
\end{proof}

\begin{definition}[Lipschitz hypothesis on the Hamiltonian]
\label{def:hamiltonian-lipschitz-condition}
\dcref{ch10:10.2-def-2}
\group{Viscosity Solutions}
\[
\begin{cases}
|H(p,x) - H(q,x)| \leq C|p - q| \\
|H(p,x) - H(p,y)| \leq C|x - y|(1 + |p|)
\end{cases}
\tag{3}
\]
for $x, y, p, q \in \R^n$ and some constant $C \geq 0$.
\end{definition}

\begin{theorem}[Uniqueness of viscosity solution]
\label{thm:uniqueness-viscosity-solution-hj}
\uses{def:viscosity-solution-finite-horizon, lem:extrema-at-terminal-time, def:hamiltonian-lipschitz-condition}
\dcref{ch10:10.2-thm-1}
\group{Viscosity Solutions}

Under assumption (3) there exists at most one viscosity solution of (1).
\end{theorem}

\begin{proof}
\sketch If two solutions differ positively, maximize their doubled-variable difference penalized by $\epsilon^{-2}(|x-y|^2+(t-s)^2)$, a small time tilt, and spatial coercive terms. Uniform continuity makes the paired points $o(\epsilon)$ apart and keeps them away from $t=0$. Applying the two viscosity inequalities and subtracting them, the Lipschitz bounds (3) force a positive tilt to be at most $o(1)$, a contradiction.
\end{proof}

\begin{remark}[Doubling of variables technique]
\label{rem:doubling-variables-technique}
\uses{thm:uniqueness-viscosity-solution-hj}
\dcref{ch10:10.2-rem-1}
\group{Viscosity Solutions}

The uniqueness argument compares two solutions by maximizing a penalized function of two space-time variables; this is the \textit{doubling of variables} technique.
\end{remark}

\section{Control Theory, Dynamic Programming}
\label{sec:hj-control-theory-dynamic-programming}

\subsection{Introduction to Control Theory}

\begin{definition}[Controlled ODE system and admissible controls]
\label{def:controlled-ode-admissible-controls}
\dcref{ch10:10.3-def-1}
\group{Control Theory}
Fix a terminal time $T>0$, a compact set $A\subset\R^m$, and a bounded Lipschitz map $\mathbf f:\R^n\times A\to\R^n$ satisfying
\[
|\mathbf f(x,a)|\leq C,
\qquad
|\mathbf f(x,a)-\mathbf f(y,a)|\leq C|x-y|.
\tag{3}
\]
For $x\in\R^n$, $0\leq t\leq T$, and a control $\alpha$, the state equation is
\[
\begin{cases}
\dot{\mathbf x}(s)=\mathbf f(\mathbf x(s),\alpha(s)) & (t<s<T),\\
\mathbf x(t)=x.
\end{cases}
\tag{1}
\]
The admissible controls are
\[
\mathcal{A} = \{\alpha : [0,T] \to A \mid \alpha(\cdot) \text{ is measurable}\}
\tag{2}
\]
and for each $\alpha\in\mathcal A$, (1) has a unique Lipschitz response $\mathbf x^{\alpha(\cdot)}$ on $[t,T]$.
\end{definition}

\begin{definition}[Running cost, terminal cost, and the cost functional]
\label{def:control-cost-functional}
\uses{def:controlled-ode-admissible-controls}
\dcref{ch10:10.3-def-2}
\group{Control Theory}

For $x\in\R^n$, $0\leq t\leq T$, and $\alpha\in\mathcal A$, define
\[
C_{x,t}[\alpha(\cdot)] := \int_t^T h(\mathbf x(s), \alpha(s)) \, ds + g(\mathbf x(T)),
\tag{4}
\]
where $\mathbf x=\mathbf x^{\alpha(\cdot)}$ solves (1). The running and terminal costs
\[
h : \R^n \times A \to \R, \quad g : \R^n \to \R
\]
are bounded and Lipschitz uniformly in $a$:
\[
\begin{cases} |h(x, a)|, |g(x)| \le C \\ |h(x, a) - h(y, a)|, |g(x) - g(y)| \le C|x - y| \end{cases} \quad (x, y \in \R^n, a \in A)
\tag{5}
\]
for some $C$.
\end{definition}

\subsection{Dynamic Programming}

\begin{definition}[Value function]
\label{def:value-function}
\uses{def:control-cost-functional}
\dcref{ch10:10.3-def-3}
\group{Control Theory}

The \textit{value function} is
\[
u(x,t) := \inf_{\alpha(\cdot) \in \mathcal{A}} C_{x,t}[\alpha(\cdot)] \quad (x \in \R^n, 0 \le t \le T).
\tag{6}
\]
\end{definition}

\begin{theorem}[Optimality conditions]
\label{thm:dynamic-programming-optimality-conditions}
\uses{def:value-function, def:control-cost-functional}
\dcref{ch10:10.3-thm-1}
\group{Control Theory}

For each $h > 0$ so small that $t + h \leq T$, we have
\[
u(x,t) = \inf_{\alpha(\cdot) \in \mathcal{A}} \left\{ \int_t^{t+h} h(\mathbf{x}(s), \alpha(s)) ds + u(\mathbf{x}(t+h), t+h) \right\},
\tag{7}
\]
where $\mathbf{x}(\cdot) = \mathbf{x}^{\alpha(\cdot)}(\cdot)$ solves the ODE (1) for the control $\alpha(\cdot)$.
\end{theorem}

\begin{proof}
\sketch Concatenate an arbitrary control on $[t,t+h]$ with an $\epsilon$-optimal continuation from the reached state to obtain one inequality in (7). Restrict an $\epsilon$-optimal control for the full interval and use the definition of the continuation value for the reverse inequality. Let $\epsilon\downarrow0$.
\end{proof}

\subsection{Hamilton-Jacobi-Bellman Equation}

\begin{lemma}[Estimates for the value function]
\label{lem:value-function-estimates}
\uses{def:value-function, def:control-cost-functional}
\dcref{ch10:10.3-lem-1}
\group{Control Theory}

There exists a constant $C$ such that
\[
|u(x,t)| \leq C,
\]
\[
|u(x,t) - u(\hat{x},\hat{t})| \leq C(|x - \hat{x}| + |t - \hat{t}|)
\]
for all $x, \hat{x} \in \R^n$, $0 \leq t, \hat{t} \leq T$.
\end{lemma}

\begin{proof}
\sketch Bounded running and terminal costs bound $u$. Comparing trajectories driven by the same control and applying Gronwall gives Lipschitz dependence on the initial state. Shifting an $\epsilon$-optimal control in time and using the bounded dynamics and Lipschitz costs gives the corresponding time estimate.
\end{proof}

\begin{theorem}[A PDE for the value function]
\label{thm:value-function-solves-hjb-equation}
\uses{lem:value-function-estimates, thm:dynamic-programming-optimality-conditions, def:value-function}
\dcref{ch10:10.3-thm-2}
\group{Control Theory}

The value function $u$ is the unique viscosity solution of this terminal-value problem for the Hamilton-Jacobi-Bellman equation:
\[
\begin{cases}
u_t + \min_{a \in A}\{f(x,a) \cdot Du + h(x,a)\} = 0 & \text{in } \R^n \times (0,T) \\
u = g & \text{on } \R^n \times \{t = T\}.
\end{cases}
\tag{22}
\]
\end{theorem}

\begin{proof}
\sketch The cost formula gives $u(x,T)=g(x)$, and the value estimates give bounded uniform continuity. At a local maximum of $u-v$, a constant control combined with (7) rules out a negative HJB residual; at a local minimum, an $o(h)$-optimal control rules out a positive residual. Thus $u$ satisfies (24), (25). The Hamiltonian (23) obeys (3), so time reversal and the uniqueness theorem give uniqueness.
\end{proof}

\begin{remark}[The Hamilton-Jacobi-Bellman Hamiltonian]
\label{rem:hjb-hamiltonian-formula}
\uses{thm:value-function-solves-hjb-equation, def:hamiltonian-lipschitz-condition}
\dcref{ch10:10.3-rem-1}
\group{Control Theory}

The Hamilton--Jacobi--Bellman PDE has the form
\[
u_t + H(Du, x) = 0 \quad \text{in } \R^n \times (0,T),
\]
for the Hamiltonian
\[
H(p,x) := \min_{a \in A}\{f(x,a) \cdot p + h(x,a)\} \quad (p, x \in \R^n).
\tag{23}
\]
The bounds (3), (5) imply that this $H$ satisfies \cref{def:hamiltonian-lipschitz-condition}.
\end{remark}

\begin{definition}[Viscosity solution of a terminal-value problem]
\label{def:viscosity-solution-terminal-value-problem}
\uses{def:viscosity-solution-cauchy-problem}
\dcref{ch10:10.3-def-4}
\group{Control Theory}

A bounded, uniformly continuous function $u$ is a \textit{viscosity solution} of (22) provided:

(i) $u = g$ on $\R^n \times \{t = T\}$,

and

(ii) for each $v \in C^\infty(\R^n \times (0,T))$

\[
\begin{cases}
\text{if } u - v \text{ has a local maximum at a point } (x_0, t_0) \in \R^n \times (0,T), \\
\text{then} \\
v_t(x_0, t_0) + H(Dv(x_0, t_0), x_0) \geq 0,
\end{cases}
\tag{24}
\]
and
\[
\begin{cases}
\text{if } u - v \text{ has a local minimum at a point } (x_0, t_0) \in \R^n \times (0,T), \\
\text{then} \\
v_t(x_0, t_0) + H(Dv(x_0, t_0), x_0) \leq 0.
\end{cases}
\tag{25}
\]

The terminal-value inequalities reverse the initial-value inequalities in \cref{def:viscosity-solution-cauchy-problem}.
\end{definition}

\begin{remark}[Correspondence with the initial-value problem via time reversal]
\label{rem:terminal-value-time-reversal}
\uses{def:viscosity-solution-terminal-value-problem, def:viscosity-solution-cauchy-problem}
\dcref{ch10:10.3-rem-2}
\group{Control Theory}

If $u$ is a viscosity solution of (22), then $w(x,t):=u(x,T-t)$ is a viscosity solution of
\[
\begin{cases}
w_t - H(Dw, x) = 0 & \text{in } \R^n \times (0,T) \\
w = g & \text{on } \R^n \times \{t = 0\}.
\end{cases}
\]
\end{remark}

\begin{remark}[Design of optimal controls]
\label{rem:design-of-optimal-controls}
\uses{thm:value-function-solves-hjb-equation}
\dcref{ch10:10.3-rem-3}
\group{Control Theory}

For an initial state $(x,t)$, a feedback trajectory satisfies
\[
\begin{cases} \dot{x}^*(s) = f(x^*(s), \alpha^*(s)) & (t < s < T) \\ x^*(t) = x, \end{cases}
\tag{41}
\]
where at each time $s$, $\alpha^*(s) \in A$ is selected so that
\[
\begin{aligned}
f(x^*(s), \alpha^*(s)) \cdot Du(x^*(s), s) + h(x^*(s), \alpha^*(s)) \\
= H(Du(x^*(s), s), x^*(s)).
\end{aligned}
\tag{42}
\]
Thus $\alpha^*(s)$ attains the minimum in (23) and is a \textit{feedback control}. Where $u$ and $\alpha^*$ are smooth, this produces a minimum-cost trajectory; at nondifferentiability points of $u$, (42) requires a generalized interpretation.
\end{remark}

\subsection{Hopf-Lax Formula Revisited}

Assume $p\mapsto H(p)$ is convex and superlinear,
\[
\lim_{|p|\to\infty}\frac{H(p)}{|p|}=+\infty,
\]
and let $g:\R^n\to\R$ be Lipschitz. Consider
\[
\begin{cases}
u_t+H(Du)=0 & \text{in }\R^n\times(0,T),\\
u=g & \text{on }\R^n\times\{t=0\}.
\end{cases}
\tag{43}
\]
The Hopf--Lax formula is
\[
u(x,t)=\min_{y\in\R^n}\left\{tL\left(\frac{x-y}{t}\right)+g(y)\right\},
\tag{44}
\]
where
\[
L(q)=\sup_{p\in\R^n}\{p\cdot q-H(p)\}.
\tag{45}
\]

\begin{theorem}[Hopf-Lax formula as viscosity solution]
\label{thm:hopf-lax-formula-viscosity-solution}
\uses{def:viscosity-solution-finite-horizon, thm:value-function-solves-hjb-equation, def:hopf-lax-formula, lem:hopf-lax-functional-identity, thm:convex-duality-hamiltonian-lagrangian, lem:hopf-lax-lipschitz-continuity}
\dcref{ch10:10.3-thm-3}
\group{Hopf-Lax Formula}

Assume in addition that $g$ is bounded. Then the unique viscosity solution of the initial-value problem (43) is given by the formula (44).
\end{theorem}

\begin{proof}
\sketch The Hopf--Lax functional identity and Lipschitz theorem give the initial data, boundedness, and regularity. If $u-v$ has a local maximum, compare the functional identity at $(x_0,t_0)$ with $(x_0-hq,t_0-h)$, divide by $h$, and optimize over $q$; convex duality yields $v_t+H(Dv)\leq0$. At a local minimum, a minimizing point in the functional identity and the opposite strict residual produce a contradiction, giving $v_t+H(Dv)\geq0$. Finite-horizon uniqueness completes the proof.
\end{proof}

\section{References}
\label{sec:hj-references}
\textbf{Viscosity solutions.} \cite{CrandallEvansLions1984}, \cite{CrandallLions1983}, \cite{CrandallEvansLions1984}.
\textbf{Control and dynamic programming.} \cite{PLLions1982}, \cite{FlemingSoner1993}, \cite{BardiCapuzzoDolcetta1997}.
